\documentclass[10pt]{article}
\usepackage[margin=0.9in]{geometry}
\usepackage{amsmath,microtype}
\title{Guide: Why a Fixed Number of Returns Is Insufficient}
\author{Collatz proof project}
\date{September 5, 2026}
\begin{document}
\maketitle
\raggedright
\section*{What Lean proves}
Choose any finite number of returns to the progression two modulo nine.
There is a positive integer whose next that many consecutive returns all
increase. They are actual consecutive visits, not a selected subsequence.
Each return takes four ordinary shortcut Collatz steps.

\section*{The construction}
Use the seeds defined by $n_0=2$ and $n_{R+1}=144n_R+299$.
The seed $n_R$ follows the odd, odd, even, odd pattern for $R$ blocks.
Each block returns to the progression and has coefficient $27/16>1$.
The intervening remainders are eight, eight, and four, so there are no
hidden visits between the displayed returns.

The proof tracks divisibility of $11n+23$ by a power of sixteen. Every
four-step block consumes one power. The recurrence supplies as many
powers as the requested finite experiment needs.

\section*{What this rules out}
We cannot prove universal descent by taking one fixed power of the return
map: every proposed fixed number of returns has an integer counterexample.
A proof still might establish descent after a number of returns depending
on the starting seed. That remains open.

The result does not construct an infinite growing orbit. A larger requested
length uses a different, larger seed. Exchanging those two quantifiers
would be an invalid inference.

\section*{How this fits the earlier results}
The earlier sixteen-step theorem applies when a first return contracts.
Here the chosen returns expand, so the theorems do not conflict. Likewise,
the equivalence between sampled contraction and nondivergence does not
promise a uniform stopping bound. No proof or disproof of Collatz is claimed.

\section*{Verification}
From \texttt{lean/}, run \texttt{lake build Collatz.GrowingReturns}.
The theorem is included in the selected axiom audit. The accompanying
paper gives the exact construction and the quantifiers that limit its scope.
\end{document}
